
\SGARefTarget{sga7.sga7i.eviii.n2-1-13}{2.1.13}. Let $\ell$ be a prime number.  For every commutative group
$E$, denote by $T_\ell(E)$ the inverse system
\begin{equation}\tag{\SGARefTarget{sga7.sga7i.eviii.n2-1-13-1}{2.1.13.1}}
 T_\ell(E)=({}_{\ell^\nu}E)_{\nu\geq0},
\end{equation}
which plainly depends functorially on $E$.  The compatibilities of
\SGARefLink{sga7.sga7i.eviii.n2-1-12}{2.1.12}, for $n,n'$ of the form $\ell^\nu,\ell^{\nu'}$ with
$\nu'\geq\nu$, show that the homomorphisms
$\alpha^{\ell^\nu}_{P,Q}$, as $\nu$ varies, define a homomorphism of
inverse systems, functorial in $P,Q$:
\begin{equation}\tag{\SGARefTarget{sga7.sga7i.eviii.n2-1-13-2}{2.1.13.2}}
 \alpha^{(\ell)}_{P,Q}:
 T_\ell(P)\otimes T_\ell(Q)
 \longrightarrow T_\ell\bigl(\uTor_1(P,Q)\bigr).
\end{equation}
Knowledge of these homomorphisms for every prime $\ell$ is plainly
equivalent to knowledge of the $\alpha^n_{P,Q}$ for every integer
$n$.  The antisymmetry relations of \SGARefLink{sga7.sga7i.eviii.n2-1-10}{2.1.10} then become the analogous
$\ell$-adic antisymmetry relations, for every $\ell$, expressed by
the anticommutativity of the diagrams
\begin{equation}\tag{\SGARefTarget{sga7.sga7i.eviii.n2-1-13-3}{2.1.13.3}}
\begin{tikzcd}[column sep=huge,row sep=huge]
 T_\ell(P)\otimes T_\ell(Q)
   \arrow[r]
   \arrow[d,"{\alpha^{(\ell)}_{P,Q}}"']
   \arrow[rd,phantom,"-" description]&
 T_\ell(Q)\otimes T_\ell(P)
   \arrow[d,"{\alpha^{(\ell)}_{Q,P}}"]\\
 T_\ell\bigl(\uTor_1(P,Q)\bigr)
   \arrow[r]&
 T_\ell\bigl(\uTor_1(Q,P)\bigr).
\end{tikzcd}
\end{equation}
Here the horizontal arrows are the symmetry arrows.

\medskip
\noindent\underline{Remark \SGARefTarget{sga7.sga7i.eviii.n2-1-14}{2.1.14}.}
The construction of the homomorphism (\SGARefLink{sga7.sga7i.eviii.n2-1-1}{2.1.1}), and the results \SGARefLink{sga7.sga7i.eviii.n2-1-10}{2.1.10}
and \SGARefLink{sga7.sga7i.eviii.n2-1-11}{2.1.11} concerning it, may be generalized to the situation in
which one is given a commutative ring $\underline A$ on $T$ and two
ideals $\underline a$ and $\underline b$ of $\underline A$.  In the
case treated above, both ideals are the ideal
$n\mathbb Z_T$ of the constant ring $\mathbb Z_T$ on $T$.  One then
defines a homomorphism, functorial in the $\underline A$-modules
$P,Q$,
\begin{equation}\tag{\SGARefTarget{sga7.sga7i.eviii.n2-1-14-1}{2.1.14.1}}
 {}_{\underline a}P\otimes{}_{\underline b}Q
 \xrightarrow{\ \alpha\ }
 \uHom_{\underline A}
 \left(
   \frac{\underline a\cap\underline b}
        {\underline a\mathbin{\cdot}\underline b},
   \uTor^{\underline A}_1(P,Q)
 \right),
\end{equation}
as the composite of the homomorphisms
\[
\begin{split}
 {}_{\underline a}P\otimes{}_{\underline b}Q
 &\longrightarrow
 \uHom(\underline A/\underline a,P)
 \otimes
 \uHom(\underline A/\underline b,Q)\\
 &\longrightarrow
 \uHom\!\left(
   \uTor_1(\underline A/\underline a,\underline A/\underline b),
   \uTor_1(P,Q)
 \right),
\end{split}
\]
using the well-known canonical isomorphism
\begin{equation}\tag{\SGARefTarget{sga7.sga7i.eviii.n2-1-14-2}{2.1.14.2}}
 \uTor_1(\underline A/\underline a,\underline A/\underline b)
 \simeq
 \frac{\underline a\cap\underline b}
      {\underline a\mathbin{\cdot}\underline b}.
\end{equation}
In fact, the sign in \SGARefLink{sga7.sga7i.eviii.n2-1-10}{2.1.10} is ``hidden'' in the choice of this
isomorphism, according as one works with the resolution
$0\longrightarrow\underline a\longrightarrow\underline A
\longrightarrow0$ of $\underline A/\underline a$, or with the
resolution $0\longrightarrow\underline b\longrightarrow\underline A
\longrightarrow0$ of $\underline A/\underline b$.  Here it becomes
the anticommutativity of the diagram
\begin{equation}\tag{\SGARefTarget{sga7.sga7i.eviii.n2-1-14-3}{2.1.14.3}}
\begin{tikzcd}[column sep=huge,row sep=huge]
 {}_{\underline a}P\otimes{}_{\underline b}Q
   \arrow[r]
   \arrow[d]
   \arrow[rd,phantom,"-1" description]&
 {}_{\underline b}Q\otimes{}_{\underline a}P
   \arrow[d]\\
 \uHom\!\left(
   \dfrac{\underline a\cap\underline b}
         {\underline a\mathbin{\cdot}\underline b},
   \uTor_1(P,Q)
 \right)
   \arrow[r]&
 \uHom\!\left(
   \dfrac{\underline a\cap\underline b}
         {\underline a\mathbin{\cdot}\underline b},
   \uTor_1(Q,P)
 \right).
\end{tikzcd}
\end{equation}
The horizontal arrows are the symmetry arrows.  When
$\underline a=\underline b$, this anticommutativity may also be
expressed by saying that the symmetry automorphism on
$\uTor_1(\underline A/\underline a,\underline A/\underline a)$ is
$-\id$.

In the special case where
$\underline a=\underline b=n\underline A$, with $n$ a section of
$\underline A$, multiplication by $n$ induces an epimorphism
\[
 \underline A/\underline a
 \longrightarrow
 \frac{\underline a\cap\underline a}
      {\underline a\mathbin{\cdot}\underline a}
 =\underline a/\underline a^2,
\]
and hence a functorial injective homomorphism, in the
$\underline A$-module $T$,
\[
 \uHom(\underline a/\underline a^2,T)
 \longrightarrow{}_{\underline a}T.
\]
Thus in this case the homomorphism (\SGARefLink{sga7.sga7i.eviii.n2-1-14-1}{2.1.14.1}) may be interpreted as a
homomorphism
\begin{equation}\tag{\SGARefTarget{sga7.sga7i.eviii.n2-1-14-4}{2.1.14.4}}
 {}_nP\otimes{}_nQ
 \longrightarrow{}_n\uTor_1(P,Q),
\end{equation}
which is closer in form to (\SGARefLink{sga7.sga7i.eviii.n2-1-1}{2.1.1}).

Continue to suppose that
$\underline a=\underline b=n\underline A$, and that $n$ is a regular
section of $\underline A$.  The spectral sequence (\SGARefLink{sga7.sga7i.eviii.n2-1-4}{2.1.4}) still makes
sense in the present setting and yields the following exact sequence
of low-degree terms, generalizing (\SGARefLink{sga7.sga7i.eviii.n2-1-5}{2.1.5}), where $Q$ is now an
$\underline A$-module annihilated by $\underline a$:
\begin{equation}\tag{\SGARefTarget{sga7.sga7i.eviii.n2-1-14-5}{2.1.14.5}}
\begin{split}
 \uTor^{\underline A}_2(P,Q)
 &\longrightarrow
 \uTor^{\underline A/n\underline A}_2(P_n,Q)
 \longrightarrow{}_nP\otimes Q
 \xrightarrow{\ \alpha\ }
 \uTor^{\underline A}_1(P,Q)\\
 &\longrightarrow
 \uTor^{\underline A/n\underline A}_1(P_n,Q)
 \longrightarrow0.
\end{split}
\end{equation}

Return now to arbitrary $\underline a$ and $\underline b$, and let
$p$ and $q$ be two sections of $\underline A$.  We then obtain the
commutative diagram below; (\SGARefLink{sga7.sga7i.eviii.n2-1-11-1}{2.1.11.1}) is essentially the special case
obtained by taking $q=1$:
\begin{equation}\tag{\SGARefTarget{sga7.sga7i.eviii.n2-1-14-6}{2.1.14.6}}
{\footnotesize
\begin{tikzcd}[row sep=large,column sep=tiny]
 &{}_{p\underline a}P\otimes{}_{q\underline b}Q
   \arrow[dl,"p\otimes q"']
   \arrow[dr,"i\otimes j"]&\\
 {}_{\underline a}P\otimes{}_{\underline b}Q
   \arrow[d,"\alpha"']&&
 {}_{pq\underline a}P\otimes{}_{pq\underline b}Q
   \arrow[d,"\alpha"]\\
 \uHom\!\left(
   \dfrac{\underline a\cap\underline b}
         {\underline a\mathbin{\cdot}\underline b},
   \uTor_1(P,Q)
 \right)
   \arrow[dr,"\mathrm{incl.}"']&&
 \uHom\!\left(
   \dfrac{pq\underline a\cap pq\underline b}
         {pq\underline a\mathbin{\cdot}pq\underline b},
   \uTor_1(P,Q)
 \right)
   \arrow[dl,"(pq)^*"]\\
 &\uHom\!\left(
   \dfrac{\underline a\cap\underline b}
         {pq\underline a\mathbin{\cdot}\underline b},
   \uTor_1(P,Q)
 \right)&
\end{tikzcd}}
\end{equation}
Here
\[
 i:{}_{p\underline a}P\longrightarrow{}_{pq\underline a}P,
 \qquad
 j:{}_{q\underline b}Q\longrightarrow{}_{pq\underline b}Q
\]
are the canonical inclusions, and $(pq)^*$ denotes the injective
homomorphism obtained, by passing to $\uHom$, from the surjective
homomorphism
\[
 \frac{\underline a\cap\underline b}
      {pq\underline a\mathbin{\cdot}\underline b}
 \longrightarrow
 \frac{pq\underline a\cap pq\underline b}
      {pq\underline a\mathbin{\cdot}pq\underline b}
\]
induced by multiplication by $pq$.  When
$\underline a=\underline b=n\underline A$, this diagram may be
rewritten in a simpler form, given in (\SGARefLink{sga7.sga7i.eviii.n2-1-11-1}{2.1.11.1}) when $q=1$.
